\section{Results}

In this section, we present the empirical results of our KADABRA and BRAVA-GNN implementations. All benchmarks were executed using our custom evaluation script, \texttt{runBenchmark.jl}, which orchestrates the invocation of both the C++ and Julia implementations across various datasets. The test instances, consisting of a mix of directed and undirected networks from the SNAP dataset, are specified in the \texttt{instances.txt} file.

Our evaluation is primarily twofold: first, we compare the baseline single-threaded performance of the reference C++ implementation of KADABRA against our Julia port. Second, we analyze the parallel scaling capabilities of our implementation as the number of threads increases.

\subsection{Implementation Comparison: C++ vs. Julia}

To establish a baseline, we first compare the execution times of the original C++ KADABRA implementation (compiled with OpenMP) against our Julia implementation. We focus on the \texttt{facebook\_combined} network, computing the betweenness centrality for all nodes (i.e., absolute error mode with $\epsilon=0.001$, $\delta=0.1$).

As shown in preliminary single-threaded experiments, the C++ implementation requires approximately 0.97 seconds to complete the sampling process and reach the confidence bound. Our Julia implementation requires approximately 1.86 seconds for the same task. The Julia port demonstrates competitive performance---being only about a factor of $1.9\times$ slower than the highly optimized C++ code on a single thread. This slight overhead is expected due to the garbage collector and dynamic dispatch, although the zero-allocation path sampling strategy implemented in our \texttt{KadabraWorkspace} heavily mitigates these issues. 

\subsection{Parallelization and Thread Scaling}

One of the primary goals of the Julia implementation was to leverage Julia's multi-threading capabilities to achieve scalable parallel performance. We executed the benchmarks varying the number of threads $T \in \{1, 2, 4, 8, 16, 24, 36, 48\}$.

By spawning worker threads that independently sample shortest paths (Phase 1) and aggregate local counts dynamically (Phase 2), we expect a near-linear speedup up to the physical core count of the machine. The check-interval for the stopping condition is intentionally scaled with the burn-in period to minimize synchronization overhead and thread contention when accessing the shared memory. 

% TODO: Insert plot or table for thread scaling results here based on runBenchmark.jl output.
While precise scaling numbers on the largest instances are still being compiled, the architecture of our implementation ensures that lock-free atomic increments during the random walk generation prevent the bottlenecks typically associated with concurrent graph traversals.